% ============================================================
%  SEKCJA 4: MIKROBIOLOGIA DROŻDŻY W MIODZIE
% ============================================================

\section{Mikrobiologia drożdży w miodzie:
obecność vs aktywność vs fałszowanie}

Mikrobiologia miodu jest dziedziną stosunkowo rzadko analizowaną
w~kontekście oceny autentyczności produktu, mimo że~drobnoustroje
(a~wśród nich drożdże osmofilne), są nieodłącznym elementem
składu każdego miodu wyprodukowanego w~warunkach naturalnych.
Niniejsza sekcja stawia i~uzasadnia centralną tezę mikrobiologiczną
artykułu: sama obecność drożdży osmofilnych w~miodzie jest zjawiskiem
biologicznie normalnym, nie~stanowi przesłanki do~stwierdzenia
niedojrzałości produktu, a~tym bardziej jego fałszowania, o~ile
aktywność wody ($a_w$) i~zawartość wody utrzymują się na~poziomach
określonych przez obowiązujące standardy.

%,----------------------------------------------------------
\subsection{Drożdże osmofilne jako naturalny składnik miodu}
%,----------------------------------------------------------

\subsubsection{Charakterystyka biologiczna drożdży osmofilnych}

Drożdże osmofilne, zwane również osmotolerancyjnymi lub
kserotolerantnymi, to~mikroorganizmy zdolne do~wzrostu i~metabolizmu
w~środowiskach o~bardzo wysokim ciśnieniu osmotycznym, tj.~przy
aktywności wody $a_w$ na~poziomie zbliżonym do~$0{,}60$--$0{,}62$
~\cite{ruegg1981,waaw2006}. W~miodzie dominującą rolę odgrywają
przede wszystkim gatunki z~rodzaju \textit{Zygosaccharomyces}:
\textit{Zygosaccharomyces rouxii} (Boutroux) Yarrow, opisany
po~raz pierwszy w~1893~roku jako drożdże izolowane z~roztworów
sacharozy, stanowi gatunek modelowy dla~osmotolerancji
grzybów. Jego minimalna aktywność wody
dla~wzrostu wynosi $a_w = 0{,}62$, co~czyni go jednym
z~najbardziej odpornych osmotolerancyjnych grzybów znanych
nauce~\cite{waaw2006}. W~warunkach poniżej tego progu
\textit{Z.~rouxii} pozostaje w~stanie kryptobiozy, komórki
zachowują żywotność i~zdolność kiełkowania, lecz nie~wykazują
aktywności metabolicznej ani wzrostu~\cite{waaw2006,ruegg1981}.

Nowsze badania przy użyciu metod kulturowalnych i~niezależnych
od~hodowli (metagenomika) rozszerzyły listę drożdży izolowanych
z~miodu o~dalsze gatunki niekonwencjonalne, w~tym \textit{Candida} spp.,
\textit{Rhodosporidiobolus ruineniae}, \textit{Clavispora lusitaniae}
i~\textit{Metschnikowia chrysoperlae}~\cite{rodriguezmachado2024}.
Badania metagenomiczne wskazują, że~\textit{Zygosaccharomyces} pozostaje
rodzajem dominującym zarówno w~miodzie pszczoły zachodniej
(\textit{Apis mellifera}), jak i~w~miodach pszczoły wschodniej (\textit{Apis cerana}),
niezależnie od~kraju pochodzenia~\cite{yeastbiodiversity2025}.

\subsubsection{Źródła kolonizacji i powszechność występowania}
Drożdże osmofilne są integralnym składnikiem ekosystemu, ula i środowiska pszczelarskiego, co nie tylko wyklucza zasadność stawiania w tym względzie ograniczeń ale i wyklucza możliwość ich eliminacji z miodu bez zastosowania ekstremalnych ingerencji technologicznych. Główne drogi kolonizacji obejmują nektar kwiatowy, w którym liczebność drożdży wynosi typowo $10^2$--$10^4$~CFU/ml w~zależności
od~gatunku rośliny, mikroklimatu i~sezonu; pyłek kwiatowy i~pierzgę, będące
nośnikami aktywnych populacji drożdży i~bakterii mlekowych;
a~także przewód pokarmowy pszczoły -- wole miodowe i~jelito
stanowią naturalny rezerwuar mikroorganizmów~\cite{alwaili2012}.
Dodatkowym źródłem są gleba, wosk i~powierzchnie ula, których
biofilm drożdżowy jest trudny do~wyeliminowania bez~środków
dezynfekcyjnych niezgodnych z~produkcją ekologiczną~\cite{alwaili2012}.

Obecność drożdży osmofilnych w~miodzie jest zatem zjawiskiem
nieuchronnym, wynikającym z~biologii procesu produkcji,
a~nie z~uchybień produkcyjnych. Potwierdza to~piśmiennictwo
z~różnych stref klimatycznych i~geograficznych, obejmujące
zarówno miody pszczoły zachodniej, jak i~wschodniej ~\cite{yeastbiodiversity2025}.
Wolumen stwierdzonych drożdży w~normalnym niefermentującym miodzie
komercyjnym wynosi zazwyczaj od~kilkudziesięciu do~kilkuset~CFU/g;
wartości powyżej 1000~CFU/g rejestrowane są w~próbkach prawidłowych
bez~oznak fermentacji, o~ile wilgotność nie~przekracza
$18\%$~\cite{analytica2020}.


%,----------------------------------------------------------
\subsection{Aktywność wody jako krytyczny parametr kontrolny}
%,----------------------------------------------------------

\subsubsection{Definicja i znaczenie aktywności wody w miodzie}

Aktywność wody ($a_w$) jest termodynamiczną miarą dostępności
wody w~produkcie spożywczym dla~reakcji chemicznych
i~aktywności mikrobiologicznej; definiowana jako stosunek
ciśnienia pary wodnej nad produktem do~ciśnienia pary
nasyconej czystej wody w~tej samej temperaturze~\cite{waaw2006}.
W~przeciwieństwie do~zawartości wody (wilgotności), wyrażanej
w~procentach masowych, $a_w$ odzwierciedla rzeczywistą
dostępność wody dla~mikroorganizmów i~jest pierwotnym
parametrem determinującym wzrost drobnoustrojów, aktywność
enzymatyczną i~nieenzymatyczne brązowienie produktu~\cite{waaw2006}.

Miód jest roztworem nadnasyconym cukrów prostych, fruktozy
($\sim 38\%$) i~glukozy ($\sim 30\%$) oraz wyższych
oligosacharydów, w~środowisku o~minimalnej zawartości wody.
Wysoka osmolarność roztworu wynikająca z~tego nasycenia
powoduje, że~cząsteczki wody są silnie związane z~cząsteczkami
cukru, a~ich aktywność termodynamiczna jest drastycznie
obniżona~\cite{zamora2006,waaw2006}. W~miodzie o~wilgotności
$\leq 20\%$ aktywność wody wynosi zazwyczaj $a_w < 0{,}60$,
choć dokładna wartość zależy od~składu cukrowego, zawartości
wody i~temperatury pomiaru~\cite{ruegg1981,gleiter2006}.
Badanie Ruegg i~Blanc~\cite{ruegg1981}, cytowane przez wszystkie
późniejsze podręczniki mikrobiologii żywności jako fundamentalne
odniesienie dla~mikrobiologii miodu, wykazało, że~dla~zakresu
wilgotności 16--21\% aktywność wody mieści się między
$a_w = 0{,}56$ a~$a_w = 0{,}63$, z~wartością $a_w = 0{,}60$
odpowiadającą wilgotności $\approx 17{,}5$--$18{,}5\%$
w~zależności od~składu cukrowego.

\subsubsection{Próg wzrostu drożdży osmofilnych a aktywność wody}

Minimalna aktywność wody wymagana do~wzrostu drożdży osmofilnych
wynosi $a_w = 0{,}61$--$0{,}62$; poniżej tego progu wzrost jest
całkowicie zahamowany, a~komórki pozostają w~stanie anabiozy
bez~aktywności metabolicznej~\cite{ruegg1981,alwaili2012}.
W~kontekście miodu oznacza~to, że~w~produkcie o~wilgotności
$\leq 18\%$ ($a_w < 0{,}60$) drożdże osmofilne są biologicznie
nieaktywne niezależnie od~ich liczebności wyrażonej w~CFU/g.

Praktyczna implikacja jest bezpośrednia: miód zawierający nawet
$10^4$~CFU/g drożdży osmofilnych, lecz o~wilgotności $17\%$
($a_w \approx 0{,}57$--$0{,}58$), nie~ulegnie fermentacji
w~żadnych warunkach przechowywania w~temperaturze pokojowej~\cite{alwaili2012}.
Dokument Analytica Laboratories formułuje tę~zasadę jako model
wieloparametryczny: fermentacja jest wysoce prawdopodobna dopiero
przy wilgotności $> 19\%$ \emph{i} liczbie drożdży $> 10$~CFU/g
jednocześnie~\cite{analytica2020}, co~\textit{a~contrario}
wyklucza dyskwalifikację produktu na~podstawie samego wskaźnika CFU/g.

Ze względów stabilności mikrobiologicznej aktywność wody jest parametrem nieporównywalnie istotniejszym niż procentowa zawartość wody. Zauważyć należy, że dla niektórych rodzajów miodów (np. miód wrzosowy) wymagania prawne dopuszczają zawartość wody na poziomach wyższych od 20\%. Wynika to z faktu, że specyficzny skład tych miodów zapewnia stabilność mikrobiologiczną przy zawartości wody przekraczającej 20\% i jednocześni miody te należy uznać za “dojrzałe” przy wartościach podanych w wymaganiach prawnych. Jeżeli badania aktywności wody wykażą, że aktywność wody w tych miodach nie przekracza progu ($a_w < 0{,}60$) to parametr ten będzie znacznie bardziej uniwersalnym parametrem dla określenia stabilności mikrobiologicznej miodu i stopnia jego dojrzałości. 

\subsubsection{Krystalizacja glukozy jako czynnik dynamicznie modyfikujący aktywność wody}

Naturalny proces krystalizacji glukozy w~miodzie jest czynnikiem
zasługującym na~szczególną uwagę w~kontekście ryzyka fermentacji,
gdyż stanowi mechanizm, przez który miód pierwotnie stabilny
mikrobiologicznie może stać się podatny na~fermentację bez~jakiejkolwiek
ingerencji zewnętrznej. Podczas krystalizacji glukoza przechodzi
z~fazy ciekłej do~fazy stałej w~postaci monowodoratu glukozy
(C$_6$H$_{12}$O$_6 \cdot$ H$_2$O). Każda cząsteczka glukozy
wiąże jedną cząsteczkę wody, efektywnie usuwając ją z~fazy
ciekłej~\cite{zamora2006}. Jednakże faza ciekła
pozostająca po~wykrystalizowaniu glukozy staje się wzbogacona
w~fruktozę i~wodę, co~powoduje \emph{wzrost} aktywności wody
w~pozostałej fazie ciekłej, mimo że~całkowita zawartość wody
w~próbce nie~ulega zmianie~\cite{zamora2006,ijfans2019}.

Zamora i~Chirife~\cite{zamora2006} w~badaniu 50~próbek
skrystalizowanych i~ponownie rozpuszczonych miodów argentyńskich
wykazali, że~krystalizacja powoduje wzrost $a_w$ fazy ciekłej
o~$0{,}03$--$0{,}04$ jednostki w~stosunku do~miodu
w~stanie płynnym. Dla~miodów o~początkowej wilgotności
$17{,}5$--$18{,}0\%$ ($a_w \approx 0{,}58$--$0{,}59$)
krystalizacja może przesunąć $a_w$ fazy ciekłej powyżej
progu $0{,}62$, aktywując tym samym populację drożdży
osmofilnych obecnych w~produkcie~\cite{zamora2006,ijfans2019}.
Zjawisko to~prowadzi do~kluczowego wniosku metodologicznego:
kontrola mikrobiologiczna miodu musi uwzględniać stan fazowy
produktu (płynny vs.~skrystalizowany), gdyż ta~sama liczba
CFU/g niesie odmienne ryzyko fermentacyjne w~zależności od~tego,
czy miód jest płynny, czy skrystalizowany. Pomiar CFU/g bez~określenia
stanu fazowego i~$a_w$ jest zatem metodologicznie niekompletny.

Mając na uwadze iż krystalizacja miodu przebiega w sposób zróżnicowany (tj. poprzez zróżnicowane formy krystaliczne co do ich struktury jak i dynamiki tworzenia) zmiany aktywności wody, a tym samych możliwość namnażania się drożdży w krystalizującym miodzie jest zmienna . Zmienność ta jest obserwowana nie tylko wglądem czasu/postępu krystalizacji ale i punktu próbobrania z opakowania . Mając na uwadze, że propagacja krystalizacji odbywa się na poziomie molekularnym to największy potencjał stężenia wody występuję w przestrzeni przy krystalicznej oraz międzykrystalicznej. Głównie na zasadzie dyfuzji fluktuacja ta dąży do wyrównania w obszarze coraz to odleglejszym od powierzchni kryształu. Zauważyć należy, że  równowaga międzyfazowa jest procesem dynamicznym, w którym kierunek przemian fazowych zdeterminowany jest przez temperaturę (zmienne warunki przechowywania). Zależnie od warunków towarzyszących  tym przemianom fazowym  (jak np. temperatura, gęstość roztworu, dostępność tlenu, napowietrzenie miodu, wolne przestrzenie tworzące się w miodzie podczas jego krystalizacji, szybkość krystalizacji wynikająca z proporcji poszczególnych cukrów itd.) potencjał  namnażania się drożdży w roztworze jest dynamiczny i  zróżnicowany w objętości miodu.  Pomiar CFU/g bez precyzyjnego określenia stanu fazowego i~$a_w$  jest zatem metodologicznie niekompletny.

Z powyższego wynika, że gdyby aktywność wody mogła stanowić o dojrzałości miodu, gdyby parametr ten został określony w wymaganiach prawnych, należy wartość aktywności wody ustalić na poziomie możliwym do osiągnięcia w czasie miodobrania, ale i winien być  zachowany (bez dalszych zabiegów technicznych) mimo różnic temperatur i przemian fazowych zachodzących na etapie jego przechowywania i dystrybucji 

%,----------------------------------------------------------
\subsection{Liczba drożdży jako wskaźnik ryzyka, nie parametr dyskwalifikujący}
%,----------------------------------------------------------

\subsubsection{Wieloparametryczny model ryzyka fermentacji}

Analiza literatury naukowej i~dokumentów praktycznych wskazuje
na~powszechne przyjęcie wieloparametrycznego modelu ryzyka
fermentacji miodu, w~którym żaden pojedynczy wskaźnik nie~jest
wystarczający do~orzeczenia o~stabilności lub niestabilności
mikrobiologicznej produktu. Ryzyko fermentacji jest funkcją
co~najmniej czterech zmiennych:

\begin{equation}
R_f = f\!\left(a_w,\; T,\; t,\; N_{CFU}\right)
\label{eq:fermentation_risk}
\end{equation}

\noindent gdzie: $R_f$, ryzyko fermentacji (prawdopodobieństwo
zapoczątkowania aktywnej fermentacji), $a_w$, aktywność wody,
$T$, temperatura przechowywania [$^\circ$C], $t$, czas
przechowywania [dni], $N_{CFU}$, liczba drożdży osmofilnych
[CFU/g].

Zgodnie z~tym modelem, drożdże mogą fermentować miód wyłącznie
gdy \emph{wszystkie} parametry sprzyjające wzrostowi są spełnione
jednocześnie: $a_w$ musi przekraczać próg $0{,}62$, temperatura
musi mieścić się w~zakresie optymalnym dla~drożdży ($10$--$27\,^\circ$C),
a~czas ekspozycji musi być wystarczający do~osiągnięcia gęstości
komórkowej inicjującej zauważalną fermentację~\cite{apinz2021}.
Dokument New Zealand Beekeeping Industry wskazuje wprost, że~miód
o~wilgotności $> 19\%$ i~liczbie drożdży $> 10$~CFU/g
jest narażony na~fermentację, co~\emph{a contrario} implikuje,
że~miód o~wilgotności $< 17\%$, nawet z~liczbą drożdży
wielokrotnie wyższą, nie~fermentuje~\cite{apinz2021}.

Tabela~\ref{tab:fermentation_risk} ilustruje macierzowy model
ryzyka fermentacji w~zależności od~wilgotności i~liczby drożdży,
zgodnie z~dostępnymi danymi empirycznymi.

\begin{table}[ht]
\begin{adjustwidth}{-\extralength}{0cm}
\centering
\caption{Szacunkowe ryzyko fermentacji miodu w~zależności
         od~wilgotności i~liczby drożdży osmofilnych
         (przechowywanie w~temperaturze $10$--$27\,^\circ$C).
         Źródło: opracowanie własne na~podstawie~\cite{analytica2020,apinz2021}.}
\label{tab:fermentation_risk}
\small
\begin{tabularx}{\fulllength}{lccc}
\hline
\multirow{2}{*}{\textbf{Wilgotność [\%]}} &
\multicolumn{3}{c}{\textbf{Liczba drożdży [CFU/g]}} \\
\cline{2-4}
 & $< 10$ & $10$--$1000$ & $> 1000$ \\
\hline
$\leq 17{,}0$ & brak ryzyka & brak ryzyka      & brak ryzyka \\
$17{,}1$--$18{,}0$ & brak ryzyka & ryzyko niskie    & ryzyko umiarkowane \\
$18{,}1$--$19{,}0$ & brak ryzyka & ryzyko umiarkowane & ryzyko wysokie \\
$19{,}1$--$20{,}0$ & ryzyko niskie & ryzyko wysokie  & fermentacja bardzo prawdopodobna \\
$> 20{,}0$    & ryzyko umiarkowane & fermentacja prawdopodobna & fermentacja nieuchronna \\
\hline
\end{tabularx}
\end{adjustwidth}
\end{table}

\subsubsection{Ograniczenia metody płytkowej (CFU/g) jako narzędzia oceny ryzyka}

Standardowa metoda oznaczania liczby drożdży i~grzybów w~żywności
opiera się na~wysiewie rozcieńczeń produktu na~podłoże selektywne
(np.~podłoże Dichloran Rose Bengal Chloramphenicol, DRBC Agar,
ISO 21527-1:2008) i~zliczaniu kolonii po~inkubacji w~temperaturze
$25\,^\circ$C przez 5~dni~\cite{intechopen_microb2017}. Metoda ta,
choć szeroko stosowana i~dobrze zwalidowana, posiada szereg
istotnych ograniczeń w~kontekście oceny ryzyka fermentacji miodu:

\begin{enumerate}
    \item \textbf{Brak rozróżnienia komórek żywotnych od~martwych:}
          metoda płytkowa zlicza wyłącznie komórki zdolne do~tworzenia
          kolonii na~podłożu stałym w~warunkach inkubacji; komórki
          żywe, lecz będące w~stanie głębokiej anabiozy
          spowodowanej niską $a_w$ miodu, mogą nie~rosnąć
          na~podłożu lub~rosnąć w~sposób atypiczny, co~prowadzi
          do~niedoszacowania prawdziwej liczby
          drożdży~\cite{intechopen_microb2017};
    \item \textbf{Brak korelacji z $a_w$ in situ:} wynik CFU/g
          jest miarą potencjalnej liczebności populacji,
          lecz nie~informuje o~tym, czy w~warunkach panujących
          w~badanym produkcie ta~populacja jest aktywna;
          dopiero połączenie CFU/g z~pomiarem $a_w$ pozwala
          ocenić rzeczywiste ryzyko fermentacji~\cite{ruegg1981};
    \item \textbf{Brak specyficzności gatunkowej:} standardowe
          podłoża selektywne nie~różnicują gatunków drożdży
          o~różnym potencjale fermentacyjnym; \textit{Metschnikowia}
          spp.\ o~niskim potencjale fermentacyjnym liczone są
          identycznie jak \textit{Z.~rouxii} o~wysokim potencjale~\cite{rodriguezmachado2024}.
\end{enumerate}

\textbf{Wniosek prawno-naukowy:} w~świetle powyższych ograniczeń
metodycznych, samo stwierdzenie liczby drożdży (CFU/g) w~miodzie, bez jednoczesnego pomiaru $a_w$, określenia stanu fazowego
produktu, temperatury przechowywania i~czasu od~zbioru,
nie~stanowi wystarczającej podstawy naukowej do~orzekania
o~ryzyku fermentacji, a~tym bardziej do~dyskwalifikacji produktu
z~obrotu handlowego. Zważywszy, że~żaden obowiązujący standard
(Codex, UE, ISO) nie~ustanawia limitu CFU/g dla~miodu
(por.~Sekcja~2), potencjalna decyzja inspektora o~dyskwalifikacji
produktu na~podstawie wyłącznie liczby CFU/g byłaby pozbawiona
zarówno podstawy naukowej, jak i~podstawy prawnej.

%,----------------------------------------------------------
\subsection{Fermentacja jako zdarzenie dynamiczne i post-harvest}
%,----------------------------------------------------------

Stanowisko Apimondia oraz znaczna część literatury traktującej
o~fałszowaniu miodu milcząco zakłada, że~fermentacja wykryta
w~miodzie jest dowodem jego niedojrzałości w~chwili zbioru.
Założenie to~jest jednak błędne z~perspektywy dynamiki
mikrobiologicznej, gdyż fermentacja jest zdarzeniem mogącym
zajść na~każdym etapie łańcucha dostaw, nie~tylko w~chwili
opuszczania ula przez produkt~\cite{apinz2021,analytica2020}.

Prawidłowo zebrany i~spełniający wszystkie kryteria normatywne
miód o~wilgotności, na przykład, $18{,}5\%$ ($a_w \approx 0{,}60$)
i~liczbie drożdży $200$~CFU/g, przechowywany przez pół roku
w~magazynie o~temperaturze $22\,^\circ$C i~wilgotności względnej
powietrza powyżej 70\%, może  miejscowo cechować się  większą aktywnością wody (krystalizacja, zmiany temepratur, dostęp powietrza itp) co~przesunie $a_w$ powyżej progu wzrostu drożdży, inicjując
fermentację~\cite{apinz2021}. 

Inspekcja graniczna, pobierająca próbki miodu importowanego
po~kilkumiesięcznym transporcie i~składowaniu, dokonuje
pomiaru w~punkcie odległym w~czasie i~przestrzeni od~momentu
zbioru; stwierdzenie fermentacji lub podwyższonej liczby
drożdży aktywnych w~tej chwili nie~pozwala, bez~analizy
dokumentacji temperatury i~wilgotności w~łańcuchu dostaw
(\textit{cold chain monitoring}), na~wiarygodne odtworzenie
stanu produktu w~chwili miodobrania (weryfikując stopień poszycia plastra) czy opuszczania pasieki~\cite{tsagkaris2021,schoder2026}.
Dlatego fermentacja nie może być dowodem na celowe fałszowanie produktu. Celowe fałszowanie produktu polegające na pozyskiwaniu miodu niedojrzałego, a~następnie zabiegi technologiczne zmierzające do jego sabilizacji mikrobiologicznej może mieć miejsce wyłącznie poprzez kontrole na etapie pozyskiwania miodu, a~skuteczne ograniczenie może zostać określenie technologii zabronionych (np. odwadnianie próżniowe). Fakt nie stosowania technologii zabronionych winem być koniecznym etapem urzędowych kontroli pasiek i~zakładów zajmujących się pozyskiwaniem i wprowadzeniem miodu do obrotu. Nie pozwala o tym wnioskawać ani liczba drożdży, ani fermentacja.  

